\documentclass[../DoAn.tex]{subfiles}
\begin{document}

\begin{center}
    \Large{\textbf{ABSTRACT}}\\
\end{center}
\vspace{1cm}

Emulation and commendation in units under the Ministry of National Defence operate under the Law on Emulation and Commendation No.\,06/2022/QH15 and its implementing decrees.
At the Military Science Academy, the personal annual track alone forms a three-tier chain: the Ministry of National Defence Certificate of Merit (BKBQP), the All-Forces Emulation Soldier title (CSTDTQ) and the Prime Minister's Certificate of Merit (BKTTCP).
Each tier has its own sliding-window cycle (2, 3 and 7 years) and several flag conditions to cross-check, making manual review on spreadsheets highly error-prone.

This thesis develops a web information system supporting the Political Department of the Academy throughout the lifecycle of seven commendation groups: personal annual, unit annual, length-of-service medals, contribution medals, commemorative medals, military flag medals and scientific achievements.
The platform covers proposal creation, automatic chain-eligibility checking, approval with the signed decision attached, audit logging and management of decision documents.
It cuts per-soldier eligibility check time from tens of minutes on spreadsheets to under one second, while preserving a complete audit trail for post-hoc verification.

The system is deployed as a web application on the Academy's internal network.
Chain-award rules sit in a central configuration table so that new awards can be added without rewriting the eligibility logic.
The deliverable supports four authorisation levels matching the military organisational structure: SuperAdmin, Admin, Manager and User.
It has been validated against the main business scenarios by combining manual black-box testing through the interface with an automated Jest suite covering the core eligibility rules.
The system is ready for pilot deployment at the Political Department.

\end{document}
